\section{Giới thiệu chung}
\label{sec:intro}

Trong thực hành khoa học dữ liệu hiện đại, mỗi dạng dữ liệu đặt ra những thách thức riêng về chất lượng mẫu, nguy cơ rò rỉ thông tin (\textit{data leakage}) và tiêu chí đánh giá mô hình. Dữ liệu bảng trong lĩnh vực tài chính thường mất cân bằng nghiêm trọng và biến đổi theo thời gian; dữ liệu văn bản báo chí tiếng Việt lại vướng vấn đề chuẩn hóa dấu thanh, phân đoạn từ ghép và các bản tin định kỳ gần giống nhau; trong khi dữ liệu ảnh đòi hỏi cân nhắc kỹ giữa độ chính xác phân loại, chi phí tài nguyên và khả năng giải thích quyết định của mạng nơ-ron sâu.

Báo cáo này trình bày kết quả nghiên cứu, trực quan hóa và thực nghiệm mô hình trên ba bài toán đại diện cho ba dạng dữ liệu trên: (1) phát hiện gian lận mở tài khoản ngân hàng trên tập dữ liệu bảng \textit{Bank Account Fraud}, (2) phân loại tin tức kinh tế -- tài chính tiếng Việt thu thập từ báo điện tử \textit{VnEconomy}, và (3) phân loại ảnh cảnh quan tự nhiên trên bộ \textit{Intel Image Classification}. Thay vì dừng ở việc huấn luyện một mô hình đơn lẻ, toàn bộ thực nghiệm được triển khai theo một nguyên tắc xuyên suốt: đặt câu hỏi định lượng ngay từ bước khám phá dữ liệu (\textit{EDA}), thiết kế phép chia tập và quy trình tiền xử lý phù hợp với bối cảnh vận hành thực tế, sau đó đối chiếu nhiều họ mô hình kèm kiểm định thống kê và phân tích lỗi.

\begin{table}[H]
  \centering
  \caption{Tóm tắt ba bài toán, quy mô dữ liệu và trọng tâm phương pháp luận.}
  \label{tab:overview_datasets}
  \begin{tabularx}{\linewidth}{l p{3.1cm} p{3.1cm} X}
    \toprule
    \textbf{Dạng dữ liệu} & \textbf{Tập dữ liệu} & \textbf{Quy mô} & \textbf{Thách thức \& Phương pháp chính} \\
    \midrule
    Bảng (\textit{Tabular}) & Bank Account Fraud (\texttt{Base.csv}) \cite{jesus2022turning} & 1.000.000 hồ sơ, 32 thuộc tính, 2 lớp & Tỷ lệ gian lận 1,10\%, dữ liệu khuyết mã hóa bằng số âm và dịch chuyển phân phối theo thời gian. Dùng \textit{Temporal Split}, đối chiếu Logistic Regression với LightGBM, đánh giá Bootstrap và công bằng theo nhóm tuổi. \\
    Văn bản (\textit{Text}) & Tin tức tài chính VnEconomy (tự thu thập) & 5.000 bài báo, 11 trường, 5 chuyên mục & Chuẩn hóa Unicode/dấu thanh tiếng Việt, cụm bài gần trùng (5,0\%) và rò rỉ theo tác giả. Dùng \textit{StratifiedGroupKFold}, TF-IDF, đối chiếu MultinomialNB, Logistic Regression và LinearSVC. \\
    Ảnh (\textit{Image}) & Intel Image Classification \cite{bansal2019intel} & 17.034 ảnh RGB (14.034 Train, 3.000 Test), 6 lớp & Giao thoa thị giác giữa các cặp cảnh quan gần giống nhau. Đối chiếu 7 cấu hình thuộc 4 hướng: \textit{From Scratch}, \textit{Linear Probe}, \textit{Fine-tuning} (ResNet-18, ViT-B/16), \textit{Zero-shot CLIP} và Grad-CAM. \\
    \bottomrule
  \end{tabularx}
\end{table}

Cấu trúc các phần tiếp theo được tổ chức tương ứng với từng bài toán: \Cref{sec:tabular} đi sâu vào bài toán dữ liệu bảng, \Cref{sec:text} phân tích ngữ liệu văn bản tiếng Việt, \Cref{sec:image} trình bày chuỗi thực nghiệm phân loại ảnh, và \Cref{sec:conclusion} tổng kết các nhận định kỹ thuật cùng giới hạn khách quan của nghiên cứu.
